\documentclass{beamer}
\usepackage[utf8]{inputenc}

\usetheme{Madrid}
\usecolortheme{default}
\usepackage{amsmath,amssymb,amsfonts,amsthm}
\usepackage{txfonts}
\usepackage{tkz-euclide}
\usepackage{listings}
\usepackage{adjustbox}
\usepackage{array}
\usepackage{tabularx}
\usepackage{gvv}
\usepackage{lmodern}
\usepackage{circuitikz}
\usepackage{tikz}
\usepackage{graphicx}

\setbeamertemplate{page number in head/foot}[totalframumber]

\usepackage{tcolorbox}
\tcbuselibrary{minted,breakable,xparse,skins}

\definecolor{bg}{gray}{0.95}
\lstset{
    language=C,
    basicstyle=\ttfamily\small,
    keywordstyle=\color{blue},
    stringstyle=\color{orange},
    commentstyle=\color{green!60!black},
    numbers=left,
    numberstyle=\tiny\color{gray},
    breaklines=true,
    showstringspaces=false,
}

%------------------------------------------------------------
\title{4.10.12}
\date{September 29,2025}
\author{EE25BTECH11011 - Navya Banavathu}

\begin{document}

\frame{\titlepage}

\begin{frame}{Question}
 The straight line $5x+4y=0$ passes through the point of intersection of the straight lines $x+2y-10 = 0$ and $2x+y+5 = 0$.
\end{frame}

\begin{frame}{Solution}
The equations of the given lines are
\begin{align}
    \myvec{1 & 2}\vec{x} &= 10 \label{eq1}\\
    \myvec{2 & 1}\vec{x} &= -5 \label{eq2}
\end{align}

Write the system as a single matrix equation:
\begin{align}
\myvec{1 & 2\\[4pt]2 & 1}\vec{x} &= \myvec{10\\[4pt]-5}
\end{align}
\end{frame}

\begin{frame}{Solution}
Solve using the inverse:
\begin{align}
\vec{x} 
&= \myvec{1 & 2\\[4pt]2 & 1}^{-1}\myvec{10\\[4pt]-5}= \frac{1}{-3}\myvec{1 & -2\\[4pt]-2 & 1}\myvec{10\\[4pt]-5} = \myvec{-\tfrac{20}{3}\\[6pt]\tfrac{25}{3}}
\end{align}

Thus, the intersection point is
\begin{align}
\vec{x} = \myvec{-\tfrac{20}{3}\\[4pt]\tfrac{25}{3}}
\end{align}
\end{frame}

\begin{frame}{Solution}

Now, check whether it lies on the line
\begin{align}
\myvec{5 & 4}\vec{x} &= 0
\end{align}

Substituting,
\begin{align}
\myvec{5 & 4}\myvec{-\tfrac{20}{3}\\[4pt]\tfrac{25}{3}}
&= \frac{1}{3}(-100+100)= 0 
\end{align}

Hence, The straight line $5x+4y=0$ passes through the point of intersection of the straight lines $x+2y-10 = 0$ and $2x+y+5 = 0$. 
\end{frame}


\begin{frame}[fragile]{ Code (C)}
\begin{block}{Part 1}
\begin{lstlisting}[language=C, basicstyle=\ttfamily\small, keywordstyle=\color{blue}, frame=single]
#include <stdio.h>

int find_intersection(double *x, double *y) {
    double a1 = 1, b1 = 2, c1 = 10;
    double a2 = 2, b2 = 1, c2 = -5;

    double det = a1*b2 - a2*b1;
    if(det == 0) return 0; // parallel

    *x = (c1*b2 - c2*b1) / det;
    *y = (a1*c2 - a2*c1) / det;

    if(5*(*x) + 4*(*y) == 0)
        return 1; // passes
    else
        return 2; // does not pass
}
\end{lstlisting}
\end{block}
\end{frame}


\begin{frame}[fragile]{Python Code:Import and Define points}
\begin{block}{Part 1}
\begin{lstlisting}[language=Python, basicstyle=\ttfamily\small, keywordstyle=\color{blue}, frame=single]
import ctypes
import matplotlib.pyplot as plt
import numpy as np

# --- Load the DLL ---
lib = ctypes.CDLL(r"C:\Users\Navya\code8.dll")

# Prepare variables to receive x and y from C
x = ctypes.c_double()
y = ctypes.c_double()

# Call the C function
result = lib.find_intersection(ctypes.byref(x), ctypes.byref(y))
xi, yi = x.value, y.value
\end{lstlisting}
\end{block}
\end{frame}


\begin{frame}[fragile]{Python Code}
\begin{block}{Part 2}
\begin{lstlisting}[language=Python, basicstyle=\ttfamily\small, keywordstyle=\color{blue}, frame=single]
print(f"Intersection point: ({xi:.2f}, {yi:.2f})")

if result == 0:
    print("Lines are parallel, no intersection.")
elif result == 1:
    print("The line 5x + 4y = 0 passes through the intersection point.")
else:
    print("The line 5x + 4y = 0 does NOT pass through the intersection point.")

\end{lstlisting}
\end{block}
\end{frame}

\begin{frame}[fragile]{Python Code: Plotting}
\begin{block}{Part 3}
\begin{lstlisting}[language=Python, basicstyle=\ttfamily\small, keywordstyle=\color{blue}, frame=single]
# Plot the lines 
x_vals = np.linspace(-10, 10, 400)

# Line 1: x + 2y = 10 => y = (10 - x)/2
y1 = (10 - x_vals)/2
# Line 2: 2x + y = -5 => y = -5 - 2*x
y2 = -5 - 2*x_vals
# Line 3: 5x + 4y = 0 => y = -5/4 * x
y3 = -5/4 * x_vals

plt.figure(figsize=(8,6))
plt.plot(x_vals, y1,label=r'$(1\ 2)\mathbf{x} = 10$',color='blue')
plt.plot(x_vals, y2,label=r'$(2\ 1)\mathbf{x} = -5$',color='green')
plt.plot(x_vals, y3,label=r'$(5\ 4)\mathbf{x} = 0$',color='red')
\end{lstlisting}
\end{block}
\end{frame}

\begin{frame}[fragile]{Python Code: Plotting}
\begin{block}{Part 4}
\begin{lstlisting}[language=Python, basicstyle=\ttfamily\small, keywordstyle=\color{blue}, frame=single]
# Mark the intersection point
plt.scatter(xi, yi, color='black', zorder=5)
plt.text(xi+0.5, yi, f'({xi:.2f},{yi:.2f})', fontsize=10)

plt.xlabel('x')
plt.ylabel('y')
plt.title('Intersection of Lines')
plt.grid(True)
plt.legend()
plt.xlim(-10, 10)
plt.ylim(-10, 10)

# Save figure as fig8.png
plt.savefig('fig8.png', dpi=300)
plt.show()
\end{lstlisting}
\end{block}
\end{frame}

\begin{frame}{Graphical Representation}
    \begin{figure}
        \centering
        \includegraphics[width=0.6\columnwidth]{../beamer/figs/fig8.png}
    \end{figure}
\end{frame}

\end{document}

